% ID: MAT-TRI-EQG-005
% Titolo: Dominio e zeri di una funzione goniometrica (3)
% Disciplina: Matematica
% Area: Trigonometria
% Argomento: Equazioni goniometriche
% Classe: Quarta secondaria di secondo grado
% Difficolta: 4
% Tipo: problema_numerico
% Tag: matematica, trigonometria, equazioni_goniometriche
% Fonte: verifica_goniometria_avanzata_anonimizzata
% Provenienza: materiale fornito dal docente; testo anonimizzato; soluzione redatta per Exergo

\begin{esercizio}
Determina il dominio e gli zeri della funzione: 
\[
f(x)=\dfrac{\cos 2x-\sin x}
{(\cos x-\sin x)(1+\sin x)}.
\]
Salvo diversa indicazione, lavora in $\mathbb R$; angoli in radianti.
\end{esercizio}

\begin{soluzione}
Il dominio è $\mathbb R\setminus\{\pi/4+k\pi,\,3\pi/2+2k\pi:k\in\mathbb Z\}$.
Il numeratore si annulla quando $1-2\sin^2x-\sin x=0$, cioè
$\sin x=1/2$ oppure $\sin x=-1$. Il secondo caso è escluso dal dominio.
Gli zeri sono $x=\pi/6+2k\pi$ e $x=5\pi/6+2k\pi$. In tutte le famiglie di soluzioni $k\in\mathbb{Z}$.
\end{soluzione}
